\documentclass{extarticle}
\usepackage[utf8]{inputenc}
\usepackage[T1]{fontenc}
\usepackage{geometry}
\usepackage{graphicx}
\usepackage{float}
\usepackage{hyperref}
\usepackage[english,indonesian]{babel}

\title{Optimasi Pemrosesan Template LaTeX Menggunakan Pendekatan Berbasis AI}
\author{Academic Support Agent}
\date{\today}

\begin{document}
\maketitle

\begin{abstract}

\end{abstract}

\chapter{PENDAHULUAN}


\chapter{TINJAUAN PUSTAKA}


% Original content (commented out):
% %! TEX program = lualatex
%\documentclass[a4paper,9pt]{extarticle}
\documentclass[a4paper]{article}
\usepackage[norsk]{babel}

\usepackage{fontspec}

\input{fonts}

\usepackage{parskip}
\usepackage{geometry}
\usepackage{multicol}
\usepackage{xcolor}
\usepackage{graphicx}
\graphicspath{{./figurer/}{./logo/}}
\usepackage{setspace}
\usepackage{booktabs}
\usepackage{amsfonts, amsmath, amsthm, amssymb}
\usepackage{lipsum}
\usepackage{verbatim}
\usepackage[colorlinks, citecolor=black, linkcolor=black, urlcolor=black]{hyperref}


\title{Faktaark -- mal SKDE}

\input{oppsett/makroer}
\input{oppsett/def.tex} % definere alle kommandoene

\newcommand{\tittel}{{\bf \textsf{Mal for faktaark}} -- {\normalfont \textsf{Undertittel}}}

\begin{document}
\newgeometry{bottom=0.5cm, right=1.5cm, left=1.5cm, top=1.5cm}
\pagestyle{empty}


\input{faktaark/ark1} % 
\input{faktaark/ark2} % 
\input{faktaark/ark3} %



\end{document}



\end{document}